\chapter{Steatorrhea}

\section*{Definition}
Steatorrhea is the presence of excess fat in the stool, defined quantitatively as fecal fat excretion exceeding 7 g per 24 hours on a 100 g/day fat diet. Clinically, it presents as pale, greasy, foul-smelling stools that float and are difficult to flush.

\section*{What Happens in Your Body}
Steatorrhea results from impaired fat digestion or absorption. Fat malabsorption occurs when pancreatic lipase is deficient (pancreatic insufficiency), intraluminal bile acid concentrations are inadequate (cholestasis, ileal disease, bacterial overgrowth), or absorptive surface area is reduced (celiac disease, short bowel syndrome). Unabsorbed triglycerides are hydrolyzed by colonic bacteria, producing fatty acids that stimulate water and electrolyte secretion, causing diarrhea.

\section*{Causes}
\begin{itemize}
  \item Pancreatic causes: Chronic pancreatitis, cystic scarring, pancreatic duct blockage (tumor, stricture), pancreatic surgical removal, Shwachman-Diamond syndrome
  \item Biliary/hepatic causes: Cholestasis (primary biliary cholangitis, primary sclerosing cholangitis), bile duct blockage (stones, strictures, malignancy), advanced liver disease reducing bile salt production
  \item Mucosal causes: Celiac disease, Tropical sprue, Whipple disease, Crohn disease, Giardiasis, Eosinophilic gastroenteritis, Autoimmune enteropathy
  \item Luminal/post-surgical: Small intestinal bacterial overgrowth (SIBO), Short bowel syndrome, Gastric bypass, Biliopancreatic diversion
  \item Other causes: Abetalipoproteinemia, Zollinger-Ellison syndrome (lipase inactivation by low pH), Medications (orlistat, acarbose)
\end{itemize}

\section*{When to See a Doctor}
See your doctor if:
\begin{itemize}
  \item Pale, greasy, foul-smelling stools persisting for more than a few days
  \item Unexplained weight loss with steatorrhea
  \item Abdominal pain or bloating accompanying fatty stools
  \item Known risk factors: pancreatitis, celiac disease, or small bowel surgery
  \item Fat-soluble vitamin deficiency symptoms: night blindness (vitamin A), easy bruising (vitamin K), bone pain (vitamin D)
\end{itemize}

\section*{Related Symptoms}
\begin{itemize}
  \item Chronic diarrhea
  \item Weight loss
  \item Abdominal bloating/distention
  \item Flatulence
  \item Foul-smelling stool
  \item Pale or clay-colored stool
  \item Muscle wasting
  \item Fat-soluble vitamin deficiencies
\end{itemize}
\textbf{Important:} A 72-hour fecal fat collection on a 100 g/day fat diet remains the gold standard, but fecal elastase-1 is a practical first-line test for pancreatic exocrine insufficiency.

\section*{Self-Care / Home Management}
\begin{itemize}
  \item Follow a low-fat diet (30--50 g/day) to reduce steatorrhea symptoms.
  \item Take medium-chain triglyceride (MCT) oil supplements as MCTs are absorbed without lipase or bile acids.
  \item Supplement fat-soluble vitamins (A, D, E, K) if malabsorption is confirmed.
  \item Avoid alcohol to prevent exacerbation of pancreatic insufficiency.
  \item Seek directed therapy if weight loss or nutritional deficiencies develop.
\end{itemize}

\section*{How Your Doctor Will Evaluate You}
\begin{itemize}
  \item 72-hour fecal fat quantification confirms steatorrhea.
  \item Fecal elastase-1 ($<$200 $\mu$g/g stool) indicates pancreatic exocrine insufficiency.
  \item Serologic testing: Tissue transglutaminase IgA (celiac disease), quantitative immunoglobulins.
  \item Breath tests for SIBO (lactulose or glucose hydrogen breath test).
  \item Imaging: CT abdomen, MRCP, or endoscopic ultrasound for pancreatic/biliary pathology.
  \item Upper endoscopy with duodenal biopsies for mucosal disease.
\end{itemize}

\section*{Treatment Options}
\begin{itemize}
  \item Pancreatic insufficiency: Pancreatic enzyme replacement therapy (PERT) with meals.
  \item Celiac disease: Strict lifelong gluten-free diet.
  \item SIBO: Antibiotics (rifaximin, metronidazole) and correction of predisposing factors.
  \item Biliary blockage: Endoscopic or surgical relief of blockage; ursodeoxycholic acid for cholestatic liver disease.
  \item Fat-soluble vitamin replacement and nutritional support as needed.
\end{itemize}

\clearpage
